% !TEX root = chapter-standalone.tex

\devel{
\section{Working with functions}
\label{sec:functions-working}

\todotext{2026-07-20, JL: explain in this section some of the main uses and perspectives one can have with functions. Especially important (in oder to interpret dependent sums and dependent products): projection/fibre picture of funcitons. For this I wonder if it might make sense to move the partitions material earlier to here for this.}




\subsection{A function is a set of choices}
Imagine specifying a function $\mapa \colon \setA \to \setB$ element by element: for each $\ela \setin \setA$ there needs to be specified a corresponding element  $\mapa(\ela) \setin \setB$. A point of view that is sometimes useful is to interpret  a function $\mapa \colon \setA \to \setB$ as encoding a set of \emph{choices}: one choice for every $\ela \setin \setA$, namely the choice of a corresponding element $\mapa(\ela)$ in the set $\setB$. 

From this point of view, it is clear, for example, that specifying a function 
\begin{equation}
\mapa \colon \makeset{\singletonel} \to \setB
\end{equation}
from a singleton set $\singleton = \makeset{\singletonel}$ is the same thing as choosing a single element of $\setB$, namely the choice of $\mapa(\singletonel)$. Thus, from this perspective, it is clear that the set
\begin{equation}
\setB^\singleton = \makeset{ \text{ functions } \singleton \to \setB \ }
\end{equation}
is in bijection with the set $\setB$ itself. 


\subsection{A function is a set of labels}

Another possible interpretation of a function $\mapa \colon \setA \to \setB$ is that $\setB$ is a set of possible \emph{labels} for elements of $\setA$, and $\mapa$ assigns each $\ela \in \setA$ a label. 

For instance, $\setA$ might be a set of foods and $\setB = \makeset{ \text{fruit}, \text{vegetable}, \text{baked}, \text{candy}}$ a set of types of foods. 

Or, $\setA$ might be a set of people in a government economic database and the elements of $\setB  = \makeset{ \text{employed}, \text{unemployed}}$ mark possible states of employment. 

\todotext{Make more examples?}

\subsection{A function is a partition}


\subsection{A function is a bundle}

The idea of a "bundle" comes from differential geometry and related fields. Intuitively speaking, a bundle consists of a ``space'' $\setB$ (imagine this to be, for example, the surface of the earth), together with, for each point $\elb \setin \setB$, a set $\subA_\elb$ of possible values of something (for example, the set of possible values of temperature that could be measured at the point $\elb$). Now, the picture to imagine is all of the sets $\subA_\elb$ hovering over the space $\setB$, forming a ``bundle over $\setB$''. 




}